\documentclass[12pt, a4paper]{article}
\usepackage[utf8]{inputenc}
\usepackage[russian]{babel}
\usepackage{amsmath, amsfonts, amssymb}
\usepackage{hyperref}
\usepackage{booktabs}
\usepackage{geometry}
\geometry{top=2cm, bottom=2cm, left=2.5cm, right=2.5cm}

\begin{document}

\section*{1. Описание проблемы и постановка задачи}

\subsection*{Короткое описание проблемы}
В задаче ректификации документов ключевым этапом является попиксельное извлечение геометрии кадра: сегментация границ страницы и предсказание карт 3D-деформации. 

Модели семейства V-JEPA обучены методом предсказания репрезентаций маскированных участков видео. Благодаря этому их пространственно-временные векторы содержат глубокое представление о трехмерной форме, структурах и геометрии объектов.

Основная проблема заключается в эффективной адаптации замороженного ViT-бэкбоуна V-JEPA под попиксельную сегментацию изогнутых страниц без высоких затрат вычислительных ресурсов на полный $fine\text{-}tuning$.

\subsection*{Постановка задачи (DID-10)}
Разработать исследовательский и инжиниринговый пайплайн адаптации V-JEPA для задачи сегментации и геометрии документов:
\begin{enumerate}
    \item \textbf{Данные:} Подготовить и адаптировать открытые наборы данных с фото страниц и эталонными сканами (\texttt{DocUNET}, \texttt{DocReal}, \texttt{DIR300}, \texttt{AnyPhotoDoc6300}).
    \item \textbf{Адаптация:} Заморозить веса бэкбоуна V-JEPA ($requires\_grad = False$) и применить методы эффективной адаптации параметров (\textbf{LoRA} / \textbf{DoRA}) к $Self\text{-}Attention$ блокам в комбинации с легковесным сегментационным декодером (FPN / U-Net Head).
    \item \textbf{Обучение и оценка:} Настроить пайплайн обучения попиксельной сегментации (передний план / $UV$-карта сдвигов) и замерить метрики качества ($mIoU$, $MS\text{-}SSIM$, $Local\ Distortion$).
\end{enumerate}

\end{document}